\chapter{循环神经网络}

序列把观测写成按时间排列的随机变量。语言模型把联合分布因式分解为条件概率
$P(x_1,\ldots,x_T)$，循环神经网络用共享参数的隐状态压缩历史，
并通过时间反向传播计算梯度。正文保留原书标题，内容以公式为主。

\section{序列模型}

记观测序列为 $x_1,\ldots,x_T$。每个时间步的合法预测只能依赖已经发生的历史，
而不能依赖未来。下面先固定条件分布，再写成可训练的特征--标签对。

\subsection{统计工具}

对任意 $t$，当前项的生成规律是条件分布
\[
x_t \sim P(x_t \mid x_{t-1},\ldots,x_1).
\]
联合分布由这些条件依次相乘得到。不同的条件独立性假设对应不同的模型类。

\subsubsection{自回归模型}

链规则给出精确分解
\[
P(x_1,\ldots,x_T)=\prod_{t=1}^{T}P(x_t \mid x_{t-1},\ldots,x_1).
\]
自回归模型直接估计右侧每个条件分布，输入是过去观测本身。
隐变量自回归则另设 $h_t=f(x_t,h_{t-1})$，再用 $P(x_t \mid h_{t-1})$ 近似条件。

\subsubsection{马尔可夫模型}

若只依赖最近一项，则得到一阶马尔可夫模型
\[
P(x_1,\ldots,x_T)=\prod_{t=1}^{T}P(x_t \mid x_{t-1}),
\qquad
P(x_1 \mid x_0)=P(x_1).
\]
此时隔一步的条件可由全概率公式消去中间变量：
\begin{align*}
P(x_{t+1} \mid x_{t-1})
&=
\frac{\sum_{x_t}P(x_{t+1},x_t,x_{t-1})}{P(x_{t-1})}
\\
&=
\sum_{x_t} P(x_{t+1} \mid x_t)\,P(x_t \mid x_{t-1}).
\end{align*}
更高阶只是把条件窗口从 $1$ 换成 $\tau$。

\subsubsection{因果关系}

时间向前时，用过去解释现在。若强行反过来因式分解，
\[
P(x_1,\ldots,x_T)=\prod_{t=T}^{1}P(x_t \mid x_{t+1},\ldots,x_T),
\]
则估计对象变成“由未来推断过去”。因果序列上正向条件通常比反向更容易稳定估计。

\subsection{训练}

用正弦加噪声生成 $x_t=\sin t+\epsilon_t$，$t=1,\ldots,1000$。
嵌入窗口长度 $\tau$ 后，监督对为
\[
y_t=x_t,
\qquad
\bm{x}_t=(x_{t-\tau},\ldots,x_{t-1}).
\]
前 $\tau$ 项没有完整历史，故丢弃；训练只用前 $600$ 个对。
于是回归映射是
\[
\hat{x}_t=f(\bm{x}_t).
\]

\subsection{预测}

单步预测把真实历史送入 $f$。多步则把已有预测再送回去。以 $\tau=4$ 为例，
\begin{align*}
\hat{x}_{605}&=f(x_{601},x_{602},x_{603},x_{604}),
\\
\hat{x}_{606}&=f(x_{602},x_{603},x_{604},\hat{x}_{605}),
\\
\hat{x}_{607}&=f(x_{603},x_{604},\hat{x}_{605},\hat{x}_{606}).
\end{align*}
当输入窗口全部换成 $\hat{x}$ 后，误差沿时间累积，这就是 $k$ 步预测。

\subsection{小结}

序列模型的核心对象是
\[
P(x_1,\ldots,x_T)=\prod_{t=1}^{T}P(x_t \mid x_{t-1},\ldots,x_1),
\]
训练必须尊重时间顺序；$k$ 步预测把估计误差迭代送入同一映射 $f$。

\subsection{练习}

核验对象是窗口长度 $\tau$、噪声为零时正弦的最小充分历史，
以及把已有 $\hat{x}$ 反馈后 $k$ 步误差如何放大。

\section{文本预处理}

文本是离散序列。预处理把它映成整数下标序列，供后面的语言模型计算
$P(x_1,\ldots,x_T)$。

\subsection{读取数据集}

语料是字符串行的列表。清洗后每行是小写字母与空格组成的序列，
标点与大小写差异被抹掉，得到原始字符流。

\subsection{词元化}

词元化是映射
\[
\operatorname{tokenize}:\ \text{字符串}
\longrightarrow
(z_1,\ldots,z_L),\qquad z_\ell\in\mathcal{Z},
\]
其中 $\mathcal{Z}$ 是词元集合，可取单词或字符。输出是词元列表的列表。

\subsection{词表}

词表是单射
\[
\operatorname{Vocab}:\ \mathcal{Z}\cup\{\texttt{<unk>}\}
\longrightarrow
\{0,1,\ldots,|\mathcal{V}|-1\}.
\]
先统计语料频率 $n(z)$，丢掉 $n(z)<m$ 的低频项，剩余按频次分配下标。
未见或已删词元映到 $\texttt{<unk>}$。可选地加入填充 $\texttt{<pad>}$、
开始 $\texttt{<bos>}$、结束 $\texttt{<eos>}$。

\subsection{整合所有功能}

把读入、词元化与查表合成
\[
\text{文本}
\longmapsto
(\operatorname{corpus},\operatorname{vocab}),
\]
其中 $\operatorname{corpus}=(x_1,\ldots,x_T)$ 且 $x_t\in\{0,\ldots,|\mathcal{V}|-1\}$。
后续字符级语言模型把整本语料当作一条长序列，而不是按行切成多条。

\subsection{小结}

预处理的终点是整数序列 $(x_1,\ldots,x_T)$ 与词表 $\operatorname{Vocab}$，
后面所有概率都定义在这套下标上。

\subsection{练习}

核验词元粒度（字符或词）以及最低频次 $m$ 如何改变 $|\mathcal{V}|$。

\section{语言模型和数据集}

语言模型的对象是词元序列的联合分布
\[
P(x_1,x_2,\ldots,x_T).
\]
估计该分布后，就能给任意句子打分，也能逐项生成后续词元。

\subsection{学习语言模型}

链规则仍然成立：
\[
P(x_1,x_2,\ldots,x_T)
=
\prod_{t=1}^{T}P(x_t \mid x_1,\ldots,x_{t-1}).
\]
例如
\begin{align*}
&P(\text{deep},\text{learning},\text{is},\text{fun})
\\
&\quad=
P(\text{deep})\,
P(\text{learning}\mid\text{deep})\,
P(\text{is}\mid\text{deep},\text{learning})
\\
&\qquad\cdot
P(\text{fun}\mid\text{deep},\text{learning},\text{is}).
\end{align*}
计数估计是
\[
\hat{P}(\text{learning}\mid\text{deep})
=
\frac{n(\text{deep, learning})}{n(\text{deep})}.
\]
低频组合用拉普拉斯平滑：
\begin{align*}
\hat{P}(x)
&=
\frac{n(x)+\epsilon_1/m}{n+\epsilon_1},
\\
\hat{P}(x'\mid x)
&=
\frac{n(x,x')+\epsilon_2\hat{P}(x')}{n(x)+\epsilon_2},
\\
\hat{P}(x''\mid x,x')
&=
\frac{n(x,x',x'')+\epsilon_3\hat{P}(x'')}{n(x,x')+\epsilon_3}.
\end{align*}

\subsection{\texorpdfstring{马尔可夫模型与$n$元语法}{马尔可夫模型与n元语法}}

截断条件窗口得到 $n$ 元语法。对长度为 $4$ 的序列，
\begin{align*}
P(x_1,x_2,x_3,x_4)
&=
P(x_1)P(x_2)P(x_3)P(x_4),
\\
P(x_1,x_2,x_3,x_4)
&=
P(x_1)P(x_2\mid x_1)P(x_3\mid x_2)P(x_4\mid x_3),
\\
P(x_1,x_2,x_3,x_4)
&=
P(x_1)P(x_2\mid x_1)P(x_3\mid x_1,x_2)P(x_4\mid x_2,x_3).
\end{align*}
三式分别是一元、二元、三元语法。

\subsection{自然语言统计}

词频近似服从齐普夫定律。按降序排列后
\[
n_i \propto \frac{1}{i^{\alpha}},
\qquad
\log n_i = -\alpha\log i + c.
\]
该幂律不仅出现在一元语法，也出现在更高阶 $n$ 元计数上，
因此长片段的经验频率会迅速变得极稀。

\subsection{读取长序列数据}

语料长度 $T$ 通常远大于一次能送入网络的步数 $n$。
把长序列切成长度为 $n$ 的子序列，标签是整体左移一项后的序列，
即用 $(x_t,\ldots,x_{t+n-1})$ 预测 $(x_{t+1},\ldots,x_{t+n})$。

\subsubsection{随机采样}

每个小批量在长序列上独立抽取起点。相邻两个批量的子序列在原文中不必相邻。
长度为 $35$、步数 $5$ 时可切出
\[
\bigl\lfloor(35-1)/5\bigr\rfloor=6
\]
个特征--标签对；批量大小为 $2$ 时得到 $3$ 个批量。

\subsubsection{顺序分区}

相邻批量的第 $i$ 条子序列在原文中也相邻。
隐状态可以跨批量延续；随机采样则每个批量必须重新初始化隐状态。
两种读法都输出形状为 $(\text{批量大小},n)$ 的 $(\bm{X},\bm{Y})$。

\subsection{小结}

语言模型估计 $P(x_1,\ldots,x_T)$；$n$ 元语法用马尔可夫截断使计数可行，
长序列则按随机采样或顺序分区切成固定长度窗口。

\subsection{练习}

核验四元语法在词表大小 $10^5$ 时要存储的计数规模，
以及随机偏移是否让窗口起点在文档上接近均匀。

\section{循环神经网络}

用一个隐状态压缩全部历史，把变长条件改成固定形状的递推：
\[
P(x_t \mid x_{t-1},\ldots,x_1)\approx P(x_t \mid h_{t-1}),
\qquad
h_t=f(x_t,h_{t-1}).
\]
参数不随 $t$ 增加。

\subsection{无隐状态的神经网络}

把一个时间步的批量输入 $\bm{X}\in\mathbb{R}^{n\times d}$ 送入普通隐层
\begin{align*}
\bm{H}&=\phi(\bm{X}\bm{W}_{xh}+\bm{b}_h),
\\
\bm{O}&=\bm{H}\bm{W}_{hq}+\bm{b}_q.
\end{align*}
不同时间步的 $\bm{H}$ 互不传递，因此无法携带 $t$ 之前的信息。

\subsection{有隐状态的循环神经网络}

加入跨时间的循环项 $\bm{W}_{hh}$：
\begin{align*}
\bm{H}_t
&=
\phi(\bm{X}_t\bm{W}_{xh}+\bm{H}_{t-1}\bm{W}_{hh}+\bm{b}_h),
\\
\bm{O}_t
&=
\bm{H}_t\bm{W}_{hq}+\bm{b}_q.
\end{align*}
其中 $\bm{X}_t\in\mathbb{R}^{n\times d}$，$\bm{H}_t\in\mathbb{R}^{n\times h}$，
$\bm{O}_t\in\mathbb{R}^{n\times q}$。同一套 $(\bm{W}_{xh},\bm{W}_{hh},\bm{W}_{hq})$
在所有 $t$ 上共享。

\subsection{基于循环神经网络的字符级语言模型}

词元取字符。批量大小为 $1$ 时，序列如 ``machine'' 的每个字符先映成向量，
再按时间送入上式。时刻 $t$ 的标签是下一个字符，输出 $\bm{O}_t$ 经 softmax
得到 $P(x_{t+1}\mid h_t)$。训练目标仍是
\[
P(x_1,\ldots,x_T)=\prod_{t=1}^{T}P(x_t \mid h_{t-1}).
\]

\subsection{困惑度（Perplexity）}

平均负对数似然
\[
\frac{1}{n}\sum_{t=1}^{n}-\log P(x_t \mid x_{t-1},\ldots,x_1)
\]
再指数化，得到困惑度
\[
\exp\Biggl(-\frac{1}{n}\sum_{t=1}^{n}\log P(x_t \mid x_{t-1},\ldots,x_1)\Biggr).
\]
它是模型在每个位置上有效分支数的几何平均：均匀乱猜约为 $|\mathcal{V}|$，
完美预测为 $1$。

\subsection{小结}

循环计算
\[
\bm{H}_t=\phi(\bm{X}_t\bm{W}_{xh}+\bm{H}_{t-1}\bm{W}_{hh}+\bm{b}_h)
\]
使 $h_t$ 携带到当前为止的历史，并用困惑度衡量 $P(x_t\mid h_{t-1})$ 的质量。

\subsection{练习}

核验下一字符的输出维必须等于 $|\mathcal{V}|$，
以及沿很长的 $t$ 反传时乘积 $\prod_j\partial h_j/\partial h_{j-1}$ 的尺度。

\section{循环神经网络的从零开始实现}

把上一节的映射在字符语料上展开：独热输入、共享循环核、裁剪后的梯度更新，
以及用困惑度衡量的生成。

\subsection{独热编码}

词表大小为 $N$ 时，下标 $i$ 映到单位向量 $\bm{e}_i\in\{0,1\}^{N}$，
$(\bm{e}_i)_j=[i=j]$。小批量时间步 $\bm{X}_t$ 的每一行都是某个 $\bm{e}_{x}$，
因而 $\bm{X}_t\in\mathbb{R}^{n\times N}$，且 $\bm{X}_t\bm{W}_{xh}$ 选出第 $x$ 行。

\subsection{初始化模型参数}

隐藏宽度 $h$ 与词表大小 $N$ 决定
\[
\bm{W}_{xh}\in\mathbb{R}^{N\times h},\
\bm{W}_{hh}\in\mathbb{R}^{h\times h},\
\bm{W}_{hq}\in\mathbb{R}^{h\times N}
\]
及偏置。输入层与输出层共用同一词表，故两边维度都是 $N$。

\subsection{循环神经网络模型}

初始隐状态 $\bm{H}_0=\bm{0}\in\mathbb{R}^{n\times h}$。一个时间步是
\begin{align*}
\bm{H}_t
&=
\tanh(\bm{X}_t\bm{W}_{xh}+\bm{H}_{t-1}\bm{W}_{hh}+\bm{b}_h),
\\
\bm{O}_t
&=
\bm{H}_t\bm{W}_{hq}+\bm{b}_q.
\end{align*}
最外层按 $t=1,\ldots,n_{\text{steps}}$ 循环，把全部 $\bm{O}_t$ 拼成输出。

\subsection{预测}

给定前缀。预热期只更新 $h_t$ 不采样；之后每步取
\[
\hat{x}_{t+1}=\operatorname*{argmax}_{x}P(x \mid h_t)
\]
或按该分布抽样，再把 $\hat{x}_{t+1}$ 作为下一输入。

\subsection{梯度裁剪}

若 $f$ 是 $L$ 利普希茨的，则
\[
\bigl|f(\bm{x})-f(\bm{y})\bigr|\le L\|\bm{x}-\bm{y}\|.
\]
沿梯度走一步后
\[
\bigl|f(\bm{x})-f(\bm{x}-\eta\bm{g})\bigr|\le L\eta\|\bm{g}\|.
\]
阈值 $\theta$ 下的裁剪
\[
\bm{g}\leftarrow\min\Bigl(1,\frac{\theta}{\|\bm{g}\|}\Bigr)\bm{g}
\]
把更新幅度限制在 $\theta$ 以内，用来抑制爆炸，但不修复消失。

\subsection{训练}

一个迭代周期内：随机采样则每批量重置 $h$；顺序分区则延续 $h$ 但截断梯度。
更新前执行上式裁剪，并用困惑度汇总不同长度序列的交叉熵。

\subsection{小结}

从零实现把
\[
(\bm{e}_{x_t},h_{t-1})\mapsto(h_t,\bm{O}_t)
\]
写成可训练循环核，预热提供可用的 $h$，裁剪保证 $\|\bm{g}\|$ 有上界。

\subsection{练习}

核验独热与嵌入层第一层等价、去掉裁剪后 $\|\bm{g}\|$ 是否发散，
以及用 $P(x)^\alpha$ 抽样如何改变生成分布。

\section{循环神经网络的简洁实现}

同一套映射改由框架中的循环层一次性求值，输出层仍需单独写出。

\subsection{定义模型}

循环层接收 $\bm{X}_{1:T}$ 与 $\bm{H}_0$，返回
\[
(\bm{H}_{1:T},\bm{H}_T)
=
\operatorname{RNN}(\bm{X}_{1:T},\bm{H}_0).
\]
这里的 $\bm{H}_{1:T}$ 还不是词表上的 logits；完整模型是
\[
\bm{O}_t=\bm{H}_t\bm{W}_{hq}+\bm{b}_q.
\]
多层时，上一层的 $\bm{H}_t^{(\ell-1)}$ 作为下一层输入。

\subsection{训练与预测}

随机初始化时前缀续写接近均匀乱猜。训练目标、裁剪与困惑度与从零实现相同，
只是 $\operatorname{RNN}$ 的逐步递推由编译后的算子完成。

\subsection{小结}

简洁实现不改变
\[
\bm{H}_t=\phi(\bm{X}_t\bm{W}_{xh}+\bm{H}_{t-1}\bm{W}_{hh}+\bm{b}_h),
\]
只把该递包装成层；词表上的 $\bm{O}_t$ 仍由单独的线性映射给出。

\subsection{练习}

核验增加隐藏层数是否仍收敛，以及用同一循环核拟合序列节中的自回归窗口 $f$。

\section{通过时间反向传播}

损失对共享参数 $w_h$ 的梯度必须沿 $h_T\to h_{T-1}\to\cdots\to h_1$ 反传，
这称为通过时间反向传播（BPTT）。

\subsection{循环神经网络的梯度分析}

一步映射与逐步损失为
\begin{align*}
h_t&=f(x_t,h_{t-1},w_h),
\\
o_t&=g(h_t,w_o),
\end{align*}
\[
L=\frac{1}{T}\sum_{t=1}^{T}l(y_t,o_t).
\]
对 $w_h$ 求偏导：
\begin{align*}
\frac{\partial L}{\partial w_h}
&=
\frac{1}{T}\sum_{t=1}^{T}
\frac{\partial l(y_t,o_t)}{\partial o_t}
\frac{\partial g(h_t,w_o)}{\partial h_t}
\frac{\partial h_t}{\partial w_h}.
\end{align*}
隐状态本身对参数的依赖是递推
\[
\frac{\partial h_t}{\partial w_h}
=
\frac{\partial f(x_t,h_{t-1},w_h)}{\partial w_h}
+
\frac{\partial f(x_t,h_{t-1},w_h)}{\partial h_{t-1}}
\frac{\partial h_{t-1}}{\partial w_h}.
\]
令
\begin{align*}
a_t&=\frac{\partial h_t}{\partial w_h},
\quad
b_t=\frac{\partial f(x_t,h_{t-1},w_h)}{\partial w_h},
\\
c_t&=\frac{\partial f(x_t,h_{t-1},w_h)}{\partial h_{t-1}},
\end{align*}
则 $a_t=b_t+c_t a_{t-1}$，展开为
\[
a_t=b_t+\sum_{i=1}^{t-1}\Biggl(\prod_{j=i+1}^{t}c_j\Biggr)b_i.
\]
代回原变量，
\begin{align*}
\frac{\partial h_t}{\partial w_h}
&=
\frac{\partial f(x_t,h_{t-1},w_h)}{\partial w_h}
\\
&\quad+
\sum_{i=1}^{t-1}
\Biggl(
\prod_{j=i+1}^{t}
\frac{\partial f(x_j,h_{j-1},w_h)}{\partial h_{j-1}}
\Biggr)
\frac{\partial f(x_i,h_{i-1},w_h)}{\partial w_h}.
\end{align*}
乘积项决定梯度爆炸或消失。

\subsubsection{完全计算}

把上式中 $i=1$ 到 $t-1$ 的和全部保留，得到精确 BPTT。
计算量随 $t$ 增长，且 $\prod c_j$ 对初值极度敏感，实践中几乎不用。

\subsubsection{截断时间步}

在 $\tau$ 步后切断求和，即把 $\partial h_{t-\tau}/\partial w_h$ 视为常数。
这给出真实梯度的近似，使模型更关注短程依赖，并对应实现里对隐状态的分离。

\subsubsection{随机截断}

用伯努利变量 $\xi_t$ 随机决定是否把梯度传向 $t-1$：
\[
z_t
=
\frac{\partial f(x_t,h_{t-1},w_h)}{\partial w_h}
+
\xi_t
\frac{\partial f(x_t,h_{t-1},w_h)}{\partial h_{t-1}}
\frac{\partial h_{t-1}}{\partial w_h}.
\]
$\mathbb{E}[z_t]$ 可与完全梯度一致，但方差更大。

\subsubsection{比较策略}

随机截断把序列切成不等长片段；规则截断切成等长窗口；完全反传保留整条链。
规则截断在偏差、方差与计算之间最常用。

\subsection{通过时间反向传播的细节}

去掉激活后的线性循环为
\begin{align*}
\bm{h}_t&=\bm{W}_{hx}\bm{x}_t+\bm{W}_{hh}\bm{h}_{t-1},
\\
\bm{o}_t&=\bm{W}_{qh}\bm{h}_t,
\end{align*}
\[
L=\frac{1}{T}\sum_{t=1}^{T}l(\bm{o}_t,y_t).
\]
输出梯度
\[
\frac{\partial L}{\partial\bm{o}_t}
=
\frac{\partial l(\bm{o}_t,y_t)}{T\cdot\partial\bm{o}_t}
\in\mathbb{R}^{q}.
\]
输出权重
\[
\frac{\partial L}{\partial\bm{W}_{qh}}
=
\sum_{t=1}^{T}
\frac{\partial L}{\partial\bm{o}_t}\bm{h}_t^{\top}.
\]
终端隐状态
\[
\frac{\partial L}{\partial\bm{h}_T}
=
\bm{W}_{qh}^{\top}\frac{\partial L}{\partial\bm{o}_T}.
\]
对 $t<T$ 反向递推
\begin{align*}
\frac{\partial L}{\partial\bm{h}_t}
&=
\bm{W}_{hh}^{\top}\frac{\partial L}{\partial\bm{h}_{t+1}}
+
\bm{W}_{qh}^{\top}\frac{\partial L}{\partial\bm{o}_t}.
\end{align*}
展开后
\[
\frac{\partial L}{\partial\bm{h}_t}
=
\sum_{i=t}^{T}
\bigl(\bm{W}_{hh}^{\top}\bigr)^{T-i}
\bm{W}_{qh}^{\top}
\frac{\partial L}{\partial\bm{o}_{T+t-i}}.
\]
于是
\begin{align*}
\frac{\partial L}{\partial\bm{W}_{hx}}
&=
\sum_{t=1}^{T}
\frac{\partial L}{\partial\bm{h}_t}\bm{x}_t^{\top},
\\
\frac{\partial L}{\partial\bm{W}_{hh}}
&=
\sum_{t=1}^{T}
\frac{\partial L}{\partial\bm{h}_t}\bm{h}_{t-1}^{\top}.
\end{align*}
$\bm{W}_{hh}^{\top}$ 的高次幂使特征值 $|\lambda|>1$ 时爆炸、$|\lambda|<1$ 时消失。

\subsection{小结}

BPTT 的核心是
\[
\frac{\partial h_t}{\partial w_h}
=
b_t+\sum_{i<t}\Biggl(\prod_{j=i+1}^{t}c_j\Biggr)b_i,
\]
截断把乘积长度限制在 $\tau$；$\bm{W}_{hh}$ 的谱半径决定梯度尺度。

\subsection{练习}

核验对称矩阵 $\bm{M}$ 满足 $\bm{M}^{k}$ 的特征值为 $\lambda_i^{k}$，
以及 $\bm{W}_{hh}^{T-t}$ 与主导特征向量对齐如何对应爆炸方向。
